% Lamb_Neha.tex --- master file, complete hand calculation.
% Compile: pdflatex, bibtex, pdflatex, pdflatex.
% \showworkedtrue  -> master with the working shown
% \showworkedfalse -> scholar version, results and sources only
\documentclass[11pt]{article}
\usepackage[T1]{fontenc}
\usepackage{lmodern}
\usepackage{amsmath,amssymb}
\usepackage{amsthm}
\usepackage{booktabs}
\usepackage{array}
\usepackage[a4paper,margin=2.2cm]{geometry}
\usepackage[hidelinks]{hyperref}

\newif\ifshowworked
\showworkedtrue
\makeatletter
\newenvironment{worked}{%
  \ifshowworked\par\smallskip\noindent\textit{Working.}\par\smallskip\else
  \setbox\z@\vbox\bgroup\fi}{\ifshowworked\par\smallskip\else\egroup\fi}
\makeatother

\newcommand{\ii}{\mathrm{i}}
\newcommand{\Ncal}{\mathcal{N}}
\newcommand{\Gcal}{\mathcal{G}}
\setlength{\parskip}{3pt}
\setlength{\parindent}{0pt}

\title{\vspace{-1.2cm}Lamb waves in a one-dimensional hexagonal piezoelectric quasicrystal plate:\\
phason dynamics, size effect and a heat equation with one lag}
\author{Dr.\ Vipin \\ Gurugram University}
\date{30 October 2026}

\begin{document}
\maketitle

\begin{abstract}
\noindent
We set up and reduce the plane-strain Lamb problem in a plate of a one-dimensional hexagonal
piezoelectric quasicrystal, thickness $2h$, free faces at $z=\pm h$, the quasiperiodic direction along the
thickness. Five fields are carried: the phonon displacements $u_x,u_z$, the phason displacement $w_z$, the
electric potential $\varphi$ and the temperature rise $\theta$. The phason balance keeps both inertia and
friction, so it is a telegraph equation and runs from a wave to a diffusion as the wavenumber falls. The
phonon stress carries the nonlocal strain gradient law on its elastic part; the phason stress, the
phonon--phason coupling, the piezoelectric and pyroelectric couplings and both thermal stresses are local.
Heat conduction is of Moore--Gibson--Thompson type in its base form and carries the thermo-phason work as
well as the thermoelastic work. The five field equations are reduced to a $5\times5$ matrix in the
thickness exponent, the degree of its characteristic polynomial is counted, its parity is decided (the
families split, for every setting of the couplings), the amplitude ratios are given through the cofactors,
the seven face conditions are derived and the dispersion matrix is assembled. One hand limit and the
phason limits close the file.
\end{abstract}

\noindent\textbf{Keywords:} guided waves; quasicrystal plate; phonon--phason coupling; nonlocal strain
gradient elasticity; Moore--Gibson--Thompson; phason relaxation.

\section{Introduction}

A quasicrystal is an alloy whose atoms sit in a pattern that is ordered but not periodic. The standard
description splits the displacement into a phonon part, which is the displacement you would see in a
photograph of the lattice, and a phason part, which describes how the aperiodic order rearranges itself
locally without moving atoms along the periodic directions. In a one-dimensional quasicrystal the
aperiodic direction is a single axis, and here that axis is the thickness direction $z$.

The elastic theory of such a medium carries the phason displacement as an independent field with its own
moduli and its own mass density. For the icosahedral case the hydrodynamics was set out by Lubensky,
Ramaswamy and Toner \cite{Lubensky1985}, and a longer account with the elastic constants is due to Levine
and co-workers \cite{Levine1985}. Both are about three-dimensional quasicrystals; they are cited for the
framework, not for hexagonal constants.

Guided waves in quasicrystal plates have been computed. Zhang, Yu, Zhang and Elmaimouni give the Lamb
problem for a one-dimensional hexagonal piezoelectric quasicrystal plate \cite{Zhang2021}; the same group
has added an integral nonlocal kernel \cite{Wang2025}, a couple-stress description \cite{Wang2022} and a
layered construction with imperfect interfaces \cite{Wang2025b}. Those four papers are the elastic
baseline. What none of them does is put a temperature field with a relaxation time into the same
dispersion problem, and none carries both phason inertia and phason friction in the plate.

The phason dynamics is not an invention of ours. Agiasofitou and Lazar write the phason momentum balance
with a friction term proportional to the phason velocity in their Eq.~(26) and the resulting telegraph
equation in their Eq.~(30), and they note that the diffusive and the propagating limits are two ends of
the same equation \cite{Agiasofitou2014}. Baggioli and Landry put the same crossover in a compact form,
$\omega^2+\ii\gamma\omega=v^2q^2$, in their Eq.~(17), with $\gamma$ the relaxation rate
\cite{Baggioli2020}; they work on the hydrodynamics of the phason mode, not on a plate. The point of a
plate is that the allowed wavenumbers are set by the thickness, so the crossover can be moved through a
measurable band by changing the plate.

Heat in a quasicrystal plate is the crowded corner. Zhang, Song, Li, Han, Yu and Yang present
generalized thermoelastic Lamb waves in one-dimensional hexagonal quasicrystal nanoplates
\cite{Zhang2026}. I could not obtain that paper: only the title, the authors and the journal came through.
It may already contain the heat operator proposed here. The first week of this project is to obtain it and
to write down which generalized thermoelastic model it uses, whether one lag or two, whether it reports
attenuation, and whether its phason is dynamic or static. Until that is answered, the novelty of the heat
equation in this paper is provisional and is stated as provisional.

The contribution claimed here is therefore narrow and it is worth stating in one sentence: the plate
problem with both phason terms, a heat equation with one relaxation time, and the nonlocal strain gradient
law on the phonon stress, with the counting and the parity done by hand.

\section{Setting and assumptions}\label{sec:setting}

The plate occupies $-h\le z\le h$ and is unbounded in $x$ and $y$. The quasiperiodic direction is $z$, and
$z$ is also the plate normal, so the same axis does both jobs: this is the geometry of the published plate
papers \cite{Zhang2021}. Plane strain applies, $u_y=0$, $\partial_y=0$, and the phason displacement has one
component, $w_z$. One plane wave travels along $x$,
\begin{equation}
\bigl(u_x,u_z,w_z,\varphi,\theta\bigr)
=\bigl(\hat u_x,\hat u_z,\hat w,\hat\varphi,\hat\theta\bigr)(z)\;
\exp\bigl[\ii(kx-\omega t)\bigr],
\qquad k=k_r+\ii k_i .
\label{eq:ansatz}
\end{equation}

The assumptions: linear phonon and phason elasticity at the reference temperature; a small temperature
rise; a linear piezoelectric coupling; no pre-stress; the phason mass density and the phason friction are
material inputs; and the two size lengths are material inputs. The phason stress is local in the main
model. Applying the nonlocal strain gradient operator to the phason stress as well is a labelled
sensitivity variant and never the main model, because no published work puts the operator on phason strain
and there is no measured phason correlation length.

\section{Notation}\label{sec:notation}

\begin{center}
\small
\begin{tabular}{p{2.1cm}p{7.7cm}p{4.2cm}}
\toprule
symbol & meaning in words & unit\\
\midrule
$2h$ & thickness; faces at $z=\pm h$ & m\\
$k=k_r+\ii k_i$ & complex wavenumber; $k_i>0$ is decay & m$^{-1}$\\
$\eta$ & exponent through the thickness, one per partial wave & m$^{-1}$\\
$u_x,u_z$ & phonon displacements & m\\
$w_z$ & phason displacement along the quasiperiodic axis & m\\
$\varphi$ & electric potential & V\\
$\theta$ & temperature rise over $T_0$ & K\\
$C_{11},C_{13},C_{33},C_{44}$ & phonon elastic moduli of the one-dimensional hexagonal class & Pa\\
$K_1,K_2$ & phason elastic moduli, on $w_{zz}$ and on $w_{zx}$ & Pa\\
$R_1,R_2,R_3$ & phonon--phason coupling moduli & Pa\\
$e_{31},e_{33},e_{15}$ & piezoelectric constants & C\,m$^{-2}$\\
$\epsilon_{11},\epsilon_{33}$ & dielectric permittivities & F\,m$^{-1}$\\
$p_3$ & pyroelectric constant & C\,m$^{-2}$K$^{-1}$\\
$\beta_1,\beta_3$ & thermal stress coefficients & Pa\,K$^{-1}$\\
$\beta^{w}$ & thermo-phason coefficient & Pa\,K$^{-1}$\\
$\rho,\rho_p$ & phonon mass density, phason mass density & kg\,m$^{-3}$\\
$\gamma$ & phason friction coefficient & kg\,m$^{-3}$s$^{-1}$\\
$\rho c_p,k_{\rm th}$ & heat capacity per volume, thermal conductivity & J\,m$^{-3}$K$^{-1}$, W\,m$^{-1}$K$^{-1}$\\
$\tau_q,\tau_T$ & lag of the heat equation and its thermal companion & s\\
$l_{\rm nl},l_g$ & nonlocal and gradient lengths & m\\
\bottomrule
\end{tabular}
\end{center}

The phason field is the one that needs a picture. Imagine a tiled floor with two tile widths that never
repeat in one direction. You can shift the pattern locally by sliding a row of tiles without moving the
floor as a whole; that slide is the phason. It costs energy, it has a small inertia, and it relaxes slowly,
which is why the phason balance carries both a second derivative and a first derivative in time.

\section{Source equations}\label{sec:source}

Kinematics, with $w_{ij}=w_{i,j}$,
\begin{equation}
\varepsilon_{ij}=\tfrac12(u_{i,j}+u_{j,i}),\qquad w_{ij}=w_{i,j},\qquad E_i=-\varphi_{,i}.
\label{eq:kin}
\end{equation}

The phonon stress, the phason stress and the electric displacement,
\begin{align}
\bigl(1-l_{\rm nl}^2\nabla^2\bigr)t_{ij}&=\sigma_{ij},&
\sigma_{ij}&=C_{ijkl}\bigl(1-l_g^2\nabla^2\bigr)\varepsilon_{kl}+R_{ijkl}w_{kl}
-e_{kij}E_k-\beta_{ij}\theta,\label{eq:B1}\\
H_{ij}&=R_{klij}\varepsilon_{kl}+K_{ijkl}w_{kl}-\beta^{w}_{ij}\theta,\label{eq:B2}\\
D_i&=e_{ikl}\varepsilon_{kl}+\epsilon_{ik}E_k+p_i\theta .\label{eq:B3}
\end{align}
Equation \eqref{eq:B1} is the nonlocal strain gradient law of Lim, Zhang and Reddy with the operators on
the elastic strain energy only; the equation numbers come through the introduction of Barretta and Marotti
de Sciarra \cite[Eqs.~(18)--(19)]{Barretta2018}, and the Lim paper itself was closed to us \cite{Lim2015}.
The phonon--phason coupling with $R_{ijkl}$ and the phason moduli $K_{ijkl}$ are the standard ones for a
one-dimensional hexagonal quasicrystal, as used for the plate problem by \cite{Zhang2021}. The
thermo-phason term $\beta^{w}_{ij}\theta$ has no source I could find: it is Duhamel's analogy extended to
the phason stress, it is ours, and the paper must say so in the sentence where it appears.

The field equations,
\begin{align}
\bigl(1-l_{\rm nl}^2\nabla^2\bigr)\rho\,\ddot u_i&=\sigma_{ij,j},\label{eq:mom}\\
H_{zj,j}&=\rho_p\ddot w_z+\gamma\dot w_z,\label{eq:phason}\\
D_{i,i}&=0,\label{eq:gauss}\\
\Bigl(1+\tau_q\frac{\partial}{\partial t}\Bigr)
\Bigl[\rho c_p\dot\theta+T_0\beta_{ij}\dot\varepsilon_{ij}
+T_0\beta^{w}w_{zz}+T_0p_i\dot E_i\Bigr]&=k_{\rm th}\nabla^2\theta .
\label{eq:heat}
\end{align}
Equation \eqref{eq:phason} with both terms is the balance of \cite[Eq.~(26)]{Agiasofitou2014}, whose
Eq.~(30) is the telegraph form. Equation \eqref{eq:heat} is the Moore--Gibson--Thompson base form with the
conductivity rate put to zero, the same operator used in the piezo-semiconductor literature
\cite{Gupta2025}; the thermo-phason work term $T_0\beta^{w}\dot w_{zz}$ is ours, and the electrocaloric term
is the textbook partner of the pyroelectric term.

\section{Calculation A: constitutive relations in components}\label{sec:calcA}

\begin{worked}
The class is one-dimensional hexagonal, $6mm$ with the quasiperiodic axis along $z$. In the two-index
notation the phonon moduli are $C_{11}=C_{22}$, $C_{13}=C_{23}$, $C_{33}$, $C_{44}=C_{55}$ and
$C_{66}=(C_{11}-C_{12})/2$, with every mixed index vanishing.
The phason moduli of this class reduce to two: $K_1$ on the phason strain along the axis, $w_{zz}$, and
$K_2$ on the transverse phason strain, $w_{zx}$.
The phonon--phason coupling reduces to three: $R_1$ and $R_2$ on the normal stresses, multiplying
$w_{zz}$, and $R_3$ on the shear, multiplying $w_{zx}$. The reason is that the only phason strain
components available are $w_{zz}$ and $w_{zx}$, because $w$ has one component and $\partial_y=0$.
Plane strain removes $y$, so of the electric tensor only $\epsilon_{11}$ and $\epsilon_{33}$ survive, and of
the piezoelectric tensor only $e_{31},e_{33},e_{15}$.
\end{worked}

With $W_1=w_{zz}=\eta\hat w$ and $W_2=w_{zx}=\ii k\hat w$, the stresses are
\begin{align}
\sigma_{xx}&=\Gcal\bigl(C_{11}\varepsilon_{xx}+C_{13}\varepsilon_{zz}\bigr)+R_1W_1-e_{31}E_z-\beta_1\theta,\\
\sigma_{zz}&=\Gcal\bigl(C_{13}\varepsilon_{xx}+C_{33}\varepsilon_{zz}\bigr)+R_2W_1-e_{33}E_z-\beta_3\theta,\\
\sigma_{xz}&=2C_{44}\Gcal\,\varepsilon_{xz}+R_3W_2-e_{15}E_x,\\
H_{zz}&=R_1\varepsilon_{xx}+R_2\varepsilon_{zz}+K_1W_1-\beta^{w}\theta,\\
H_{zx}&=R_3\,2\varepsilon_{xz}+K_2W_2,
\end{align}
and the displacements,
\begin{equation}
D_x=e_{15}\,2\varepsilon_{xz}+\epsilon_{11}E_x,\qquad
D_z=e_{31}\varepsilon_{xx}+e_{33}\varepsilon_{zz}+\epsilon_{33}E_z+p_3\theta .
\end{equation}
The size-effect factor $\Gcal$ multiplies only the phonon elastic strain. The $R$ terms are local: that is
the modelling choice stated in Section~\ref{sec:setting}.

\section{Calculation B and C: the five equations and the matrix}\label{sec:calcBC}

With the plane wave in, $\nabla^2\to L=\eta^2-k^2$ and
\begin{equation}
\Ncal(\eta)=1+l_{\rm nl}^2(k^2-\eta^2),\qquad \Gcal(\eta)=1+l_g^2(k^2-\eta^2).
\label{eq:NG}
\end{equation}
Collecting the five unknown amplitudes in $\mathbf a=(\hat u_x,\hat u_z,\hat w,\hat\varphi,\hat\theta)^T$
gives $M(\eta)\mathbf a=\mathbf 0$. The entries follow.

\subsection*{Row 1: phonon momentum in $x$}

\begin{worked}
Strains: $\varepsilon_{xx}=\ii kX$, $\varepsilon_{zz}=\eta Z$, $2\varepsilon_{xz}=\eta X+\ii kZ$,
$W_1=\eta W$, $W_2=\ii kW$, $E_x=-\ii k\Phi$, $E_z=-\eta\Phi$.
$\partial_x\sigma_{xx}$ brings $\ii k$: from $u_x$, $-\Gcal C_{11}k^2$; from $u_z$, $\ii k\eta\Gcal C_{13}$;
from $w$, $\ii k\eta R_1$; from $\varphi$, $\ii k\eta e_{31}$; from $\theta$, $-\ii k\beta_1$.
$\partial_z\sigma_{xz}$ brings $\eta$: from $u_x$, $+\Gcal C_{44}\eta^2$; from $u_z$,
$\ii k\eta\Gcal C_{44}$; from $w$, $\ii k\eta R_3$; from $\varphi$, $\ii k\eta e_{15}$.
Adding and moving the inertia $-\Ncal\rho\omega^2u_x$ across gives the row.
\end{worked}

\begin{equation}
M_{11}=-\Ncal\rho\omega^2+\Gcal(C_{11}k^2-C_{44}\eta^2),\quad
M_{12}=-\ii k\eta\Gcal(C_{13}+C_{44}),\quad
M_{13}=-\ii k\eta(R_1+R_3),\quad
M_{14}=-\ii k\eta(e_{31}+e_{15}),\quad
M_{15}=\ii k\beta_1 .
\end{equation}

\subsection*{Row 2: phonon momentum in $z$}

\begin{worked}
$\partial_x\sigma_{xz}$ brings $\ii k$: $u_x$ gives $\ii k\eta\Gcal C_{44}$, $u_z$ gives
$-\Gcal C_{44}k^2$, $w$ gives $\ii k\cdot\ii kR_3=-R_3k^2$, $\varphi$ gives $-e_{15}k^2$.
$\partial_z\sigma_{zz}$ brings $\eta$: $u_x$ gives $\ii k\eta\Gcal C_{13}$, $u_z$ gives $\Gcal C_{33}\eta^2$,
$w$ gives $R_2\eta^2$, $\varphi$ gives $e_{33}\eta^2$, $\theta$ gives $-\eta\beta_3$.
Collect, and move the inertia term across.
\end{worked}

\begin{equation}
M_{21}=-\ii k\eta\Gcal(C_{13}+C_{44}),\quad
M_{22}=-\Ncal\rho\omega^2+\Gcal(C_{44}k^2-C_{33}\eta^2),\quad
M_{23}=R_3k^2-R_2\eta^2,\quad
M_{24}=e_{15}k^2-e_{33}\eta^2,\quad
M_{25}=\eta\beta_3 .
\end{equation}

\subsection*{Row 3: the phason balance}

\begin{worked}
The balance is $H_{zx,x}+H_{zz,z}=\rho_p\ddot w+\gamma\dot w$. With
$H_{zz}=R_1\ii kX+R_2\eta Z+K_1\eta W-\beta^{w}\Theta$ and
$H_{zx}=R_3(\eta X+\ii kZ)+K_2\ii kW$,
$\partial_zH_{zz}$ gives $\ii k\eta R_1$ on $X$, $R_2\eta^2$ on $Z$, $K_1\eta^2$ on $W$ and
$-\eta\beta^{w}$ on $\Theta$;
$\partial_xH_{zx}$ gives $\ii k\eta R_3$ on $X$, $-R_3k^2$ on $Z$ and $-K_2k^2$ on $W$.
The right side is $-\rho_p\omega^2W-\ii\gamma\omega W$; moving it across adds $\rho_p\omega^2+\ii\gamma\omega$
to the $W$ entry.
\end{worked}

\begin{equation}
M_{31}=\ii k\eta(R_1+R_3),\quad
M_{32}=R_2\eta^2-R_3k^2,\quad
M_{33}=K_1\eta^2-K_2k^2+\rho_p\omega^2+\ii\gamma\omega,\quad
M_{35}=\eta\beta^{w}.
\end{equation}
That entry $M_{33}$ is the whole point of the paper: divide it by the density and it is
$v_p^2\eta^2+\ii\gamma\omega/\rho_p$ against $K_2k^2$, so at large wavenumber the phason propagates and at
small wavenumber it diffuses.

\subsection*{Row 4: Gauss law in a dielectric}

\begin{worked}
$\partial_xD_x$ brings $\ii k$: $u_x$ gives $\ii k\eta e_{15}$, $u_z$ gives $-e_{15}k^2$, $\varphi$ gives
$-\epsilon_{11}k^2$. $\partial_zD_z$ brings $\eta$: $u_x$ gives $\ii k\eta e_{31}$, $u_z$ gives
$e_{33}\eta^2$, $\varphi$ gives $\epsilon_{33}\eta^2$, $\theta$ gives $\eta p_3$. There is no free charge in
a dielectric, so the sum is simply set to zero.
\end{worked}

\begin{equation}
M_{41}=\ii k\eta(e_{15}+e_{31}),\quad
M_{42}=e_{33}\eta^2-e_{15}k^2,\quad
M_{44}=\epsilon_{33}\eta^2-\epsilon_{11}k^2,\quad
M_{45}=\eta p_3 .
\end{equation}

\subsection*{Row 5: the heat equation}

\begin{worked}
Put $f=1-\ii\omega\tau_q$. The bracket of \eqref{eq:heat} has four pieces: $-\ii\omega\rho c_p\Theta$;
the thermoelastic work $T_0[\beta_1\dot\varepsilon_{xx}+\beta_3\dot\varepsilon_{zz}]
=T_0\omega k\beta_1X-\ii\omega T_0\beta_3\eta Z$; the thermo-phason work
$T_0\beta^{w}\dot w_{zz}=-\ii\omega T_0\beta^{w}\eta W$; and the electrocaloric term
$T_0p_3\dot E_z=\ii\omega T_0p_3\eta\Phi$. Multiply by $f$ and subtract $k_{\rm th}L\Theta$.
\end{worked}

\begin{equation}
M_{51}=fT_0\omega k\beta_1,\quad
M_{52}=-\ii\omega fT_0\beta_3\eta,\quad
M_{53}=-\ii\omega fT_0\beta^{w}\eta,\quad
M_{54}=\ii\omega fT_0p_3\eta,\quad
M_{55}=-\ii\omega f\rho c_p-k_{\rm th}L .
\end{equation}

\section{Calculation D: degree and parity}\label{sec:calcD}

\begin{worked}
Count the highest power of $\eta$ in each entry. Row 1: four (from $\Gcal C_{44}\eta^2$), three, three,
three, one. Row 2: three, four, two, two, one. Row 3: one, two, two, zero, one. Row 4: one, two, zero,
two, one. Row 5: zero, one, one, one, two.
The diagonal product is $4+4+2+2+2=14$ and no permutation beats it, because giving up a four in row 1 or
row 2 costs two and the best available replacement gains two. So the degree is $14$ with $l_g\ne0$. With
$l_g=0$ both fours drop to twos and the degree is $10$; the nonlocal length sits on the inertia and adds
nothing to the order.
For the parity, test whether $M_{ij}(-\eta)=\sigma_i\tau_jM_{ij}(\eta)$ admits a solution in signs.
Take $\tau_1=+1$: $M_{12}$ odd gives $\sigma_1\tau_2=-1$, $M_{13}$ odd gives $\sigma_1\tau_3=-1$,
$M_{14}$ odd gives $\sigma_1\tau_4=-1$, $M_{15}$ even gives $\sigma_1\tau_5=+1$. Choose $\sigma_1=+1$,
so $\tau_2=\tau_3=\tau_4=-1$ and $\tau_5=+1$.
Row 2: $M_{21}$ odd gives $\sigma_2=-1$; $M_{23},M_{24}$ even give $\sigma_2\tau_3=+1$ and
$\sigma_2\tau_4=+1$, true; $M_{25}$ odd gives $\sigma_2\tau_5=-1$, true.
Row 3: $M_{31}$ odd gives $\sigma_3=-1$; $M_{32}$ even and $M_{33}$ even give $\sigma_3\tau_2=+1$,
true; $M_{35}$ odd gives $\sigma_3\tau_5=-1$, true.
Row 4: $M_{41}$ odd gives $\sigma_4=-1$; $M_{45}$ odd gives $\sigma_4\tau_5=-1$, true; $M_{44}$ even,
true.
Row 5: $M_{52},M_{53},M_{54}$ odd give $\sigma_5\tau_j=-1$ for $j=2,3,4$, so with
$\tau_j=-1$ we need $\sigma_5=+1$; $M_{51}$ even then gives $+1$, true, and $M_{55}$ even, true.
The matching closes for every setting of the couplings.
\end{worked}

So here the answer is the classical one and it should be said plainly: $\det M$ is even in $\eta$ for all
values of the couplings, the fourteen roots fall into seven pairs $\pm\eta_m$, and the plate modes split
into a symmetric and an antisymmetric family. The split can be exhibited by writing the two smaller
matrices, one for each family, each of order seven; a numerical search that does not show the pairing has
gone wrong in the search, not found new physics. The contrast with the companion semiconductor paper is worth
one sentence in the article: there the pyroelectric charge and the carrier field break the mirror, and here
there is no carrier field to break it.

\section{Calculation E: amplitude ratios and the solution}\label{sec:calcE}

\begin{worked}
Take the seven roots with $\operatorname{Re}\eta\ge0$ for each family, or all fourteen when the full matrix
is used. For each root the amplitude vector is a column of the cofactor matrix
$C_{ij}=(-1)^{i+j}\mathcal M_{ji}$, where $\mathcal M_{ji}$ is the minor with row $j$ and column $i$
deleted, because $M(\eta_m)C=0$ whenever $\det M(\eta_m)=0$.
A practical choice: delete the row with the largest norm at that root and use the corresponding column, so
that the vector is not the zero vector and is well scaled. The arbitrary scalar is absorbed in $A_m$.
\end{worked}

\begin{equation}
\begin{pmatrix}u_x\\u_z\\w\\\varphi\\\theta\end{pmatrix}
=\sum_{m=1}^{14}A_m
\begin{pmatrix}X_m\\Z_m\\W_m\\\Phi_m\\\Theta_m\end{pmatrix}
e^{\eta_m z}e^{\ii(kx-\omega t)} ,
\end{equation}
with the amplitudes normalised so that the largest component is one, and fourteen constants fixed by the
face conditions.

\section{Calculation F: the seven face conditions}\label{sec:calcF}

\begin{worked}
The tractions and the double tractions are the same construction as in the companion paper: the free-face
traction of a strain-gradient body carries the surface-gradient term, and the double traction is the
quantity the gradient term does work on, both as in \cite[Sec.~2, Eq.~(14)]{Charal2019}.
For a root $\eta_m$ the seven rows are the normal traction, the shear traction, the two double tractions,
the normal electric displacement, the normal phason stress and the normal heat flux, in that order.
The phason condition is the one that has no counterpart in an elastic plate: the phason stress normal to
the face vanishes, $H_{zz}=0$, which is the condition used for a free surface in the quasicrystal
literature of \cite{Zhang2021}.
\end{worked}

\begin{align}
b_1&=\frac{1}{\Ncal_m}\Bigl[\hat\sigma_{zz}-l_g^2\ii k\eta\,C_{44}(\eta X_m+\ii kZ_m)\Bigr],&
b_2&=\frac{1}{\Ncal_m}\Bigl[\hat\sigma_{xz}-l_g^2\ii k\eta\,(\ii kC_{11}X_m+\eta C_{13}Z_m)\Bigr],\\
b_3&=\frac{l_g^2}{\Ncal_m}\eta\,C_{44}(\eta X_m+\ii kZ_m),&
b_4&=\frac{l_g^2}{\Ncal_m}\eta\,(\ii kC_{13}X_m+\eta C_{33}Z_m),\\
b_5&=\ii ke_{31}X_m+\eta e_{33}Z_m-\eta\epsilon_{33}\Phi_m+p_3\Theta_m,&
b_6&=R_1\ii kX_m+R_2\eta Z_m+K_1\eta W_m-\beta^{w}\Theta_m,\\
b_7&=\eta\Theta_m .&&
\end{align}
Rows $b_1$ and $b_2$ use the stresses of Calculation A at the root, with the $R$ terms inside them.

\section{Calculation G: the dispersion matrix}\label{sec:calcG}

\begin{equation}
\mathbf B(k,\omega)=
\begin{pmatrix}b_r(\eta_m)e^{\eta_mh}\\ b_r(\eta_m)e^{-\eta_mh}\end{pmatrix}_{r=1\dots7,\;m=1\dots14},
\qquad \det\mathbf B=0 ,
\end{equation}
solved for complex $k$ at real $\omega$. Because the parity holds, the same problem can be posed twice with
seven roots and seven rows, once for each family; the two routes must agree, and checking that they do is
one of the self-tests. When $l_g=0$ the rows $b_3,b_4$ drop and ten roots are used.

\section{Calculation H: hand limits}\label{sec:calcH}

\subsection*{Phason off}

\begin{worked}
Set $R_1=R_2=R_3=0$, $K_1=K_2=0$ and drop the phason row and column: what is left is the piezoelectric
thermoelastic plate, with a $4\times4$ matrix and the classical families. That is the reduction that must be
checked first, because it uses only published physics.
\end{worked}

\subsection*{The two phason limits}

\begin{worked}
Look at $M_{33}=K_1\eta^2-K_2k^2+\rho_p\omega^2+\ii\gamma\omega$. Divide by $\rho_p$ and put
$\omega=0$: it becomes $K_1\eta^2/\rho_p-K_2k^2/\rho_p$, which is the static balance and gives the phason
diffusion length. Keep $\gamma=0$ and $\omega\ne0$: the balance is $K_1\eta^2+\rho_p\omega^2=K_2k^2$, a wave
equation with speed $\sqrt{K_1/\rho_p}$ in the thickness direction. Keep $\rho_p=0$: the balance is
$K_1\eta^2+\ii\gamma\omega=K_2k^2$, a diffusion equation with rate $K_1/\gamma$. So the plate carries a
propagating phason branch at large wavenumber and a diffusive one at small wavenumber, with the crossover
at $\omega\sim\gamma/\rho_p$, exactly the crossover of \cite[Eq.~(17)]{Baggioli2020} written with their
$\gamma$ as the relaxation rate.
\end{worked}

The result to check is that switching the phason off returns the piezoelectric thermoelastic plate, and
that the phason branch changes character at $\omega\sim\gamma/\rho_p$ with the two limits above.

\section{Self-test list}\label{sec:selftest}

\begin{enumerate}
\item All couplings off and both size lengths zero: the classical Rayleigh--Lamb plate, both families.
\item Phason off: the piezoelectric thermoelastic plate of the companion paper, same branch count.
\item $\gamma=0$ with $\rho_p\ne0$: a propagating phason branch; $\rho_p=0$ with $\gamma\ne0$: a diffusive
one; both with the same crossover.
\item The parity statement of Calculation D, checked twice: from the roots (do they pair?) and from the two
smaller family matrices (do their determinants vanish at the same $k$?).
\item $l_g\to0$ and $l_{\rm nl}\to0$ separately, continuously, with the gradient length stiffening and the
nonlocal length softening.
\item The count of rows in $\mathbf B$ against the count of roots, in both settings of $l_g$.
\item Dimensions of rows 1--3 against each other; $R$ and $K$ must enter with the same unit.
\end{enumerate}

\section{Work plan}\label{sec:plan}

Week one is the reading that decides the novelty: obtain \cite{Zhang2026} and write down its heat operator,
its phason treatment and whether it reports attenuation. Weeks two to three: the matrix by hand from the
field equations, the root count, the parity, and the two family matrices. Weeks four to five: the phason
limits, the size effect and the branch ordering against the published elastic plate \cite{Zhang2021}.
Weeks six to eight: dispersion and attenuation curves against $f(2h)$, for two values of the phason
friction. Weeks nine to ten: the crossover map, one figure. Weeks eleven to twelve: writing and the
citation check.

\bibliographystyle{unsrt}
\bibliography{Lamb_Neha}

\end{document}
